\section{Introduction}

Large language models (LLMs) have exhibited impressive capabilities in solving diverse tasks through natural language, leading to numerous successful chat assistant applications \citep{openai2022chatgpt, microsoft2023copilot}. Nevertheless, LLMs face limitations on tasks relying heavily on personal knowledge accumulated through long-term user-AI interactions, such as psychological counseling or secretarial duties \citep{zhong2024memorybank}. Failing to incorporate user background and preferences into responses can diminish the response's accuracy as well as user satisfaction. To personalize LLM-based assistants, long-term memory, the ability to memorize, recall, and reason with a long interaction history, is indispensable. Recently, several commercial \citep{openai_memory_2024, coze_memory_2024} and open-source assistant systems with memory \citep{zhong2024memorybank, zhang2024cognitivekernelopensourceagent} have been introduced. These systems leverage techniques like compressing, indexing, and retrieving from chat histories to generate more accurate and personalized responses. 


Despite these advances, there has been limited progress in holistically evaluating the memory capability in long-term interactions. While several benchmarks evaluate LLMs on understanding long chat histories \citep{xu-etal-2022-beyond, xu-etal-2022-long, zhong2024memorybank, maharana-etal-2024-evaluating, du-etal-2024-perltqa, kim2024dialsimrealtimesimulatorevaluating}, they have two major shortcomings. First, they do not accurately reflect user-AI interactions: many focus solely on human-human conversations \citep{xu-etal-2022-beyond, maharana-etal-2024-evaluating, kim2024dialsimrealtimesimulatorevaluating}, while others omit task-oriented dialogues, which represent a significant portion of chat assistant usage and challenge memorization with the long-context inputs and long-form responses. Their interactive histories also typically have a non-configurable length spanning only a few thousand tokens, limiting the difficulty as current systems continue to improve. Second, current benchmarks' questions only offer a limited coverage of the memory abilities required in dynamic long-term interactions. For instance, MemoryBank \citep{zhong2024memorybank} and PerLTQA \citep{du-etal-2024-perltqa} insufficiently evaluate the ability to synthesize information across numerous sessions or to reason with temporal metadata or time references. All long-term memory benchmarks including recent ones such as LoCoMo \citep{maharana-etal-2024-evaluating} also fail to evaluate recall of information provided by the assistant or reasoning with updated user information.  

We introduce \BENCHMARK, a comprehensive benchmark for assessing the long-term memory capabilities of chat assistants. \BENCHMARK{} consists of 500 manually created questions to test five core memory abilities: information extraction, multi-session reasoning, temporal reasoning, knowledge updates, and abstention. Each question requires recalling information hidden within one or more task-oriented dialogues between a user and an assistant. Inspired by the ``needle-in-a-haystack" test \citep{kamradt2023needle}, we design a pipeline to compile a coherent and length-configurable chat history for each question. A chat system, then, is required to parse the dynamic interactions online for memorization, and answer the question after all the interaction sessions. While the length of the history is freely extensible, we provide two standard settings for consistent comparison: \BENCHMARK\textsubscript{\textsc{S}} with approximately 115k tokens per problem and \BENCHMARK\textsubscript{\textsc{M}} with 500 sessions (around 1.5 million tokens). Preliminary evaluations highlight the difficulty of \BENCHMARK, as long-context LLMs show a 30\%$\sim$60\% performance drop on \BENCHMARK\textsubscript{\textsc{S}}, and manual evaluations reveal that state-of-the-art commercial systems only achieved 30\%$\sim$70\% accuracy in a setting much simpler than \BENCHMARK\textsubscript{\textsc{S}} (\cref{section-benchmark-challenging}).

Finally, we present a unified view for memory-augmented chat assistants. Leveraging \BENCHMARK, we comprehensively analyze memory design choices across three key execution stages—indexing, retrieval, and reading—and four control points: value, key, query, and reading strategy. Experimental insights identify several effective memory designs: 

\begin{itemize}
    \item (\Cref{section-value-results}) Instead of sessions, \textit{round} is the more optimal granularity for storing and utilizing the interactive history. While further compression into individual user facts harms overall performance due to information loss, it improves the multi-session reasoning accuracy. 
    \item (\Cref{section-key-results}) While using a flat index with the memory values themselves as the keys is a strong baseline, further \textit{expanding} the keys with extracted user facts improves both memory recall {(9.4\% higher recall@$k$) and downstream question answering (5.4\% higher accuracy}). 
    \item (\Cref{section-query-results}) Naive time-agnostic memory designs perform poorly on temporal reasoning questions. We propose a simple indexing and query expansion strategy to explicitly associate timestamps with facts and narrow down the search range, improving the memory recall for temporal reasoning {by 6.8\%$\sim$11.3\% when a strong LLM is employed for query expansion}.  
    \item (\Cref{section-reading-results}) Even with perfect memory recall, accurately utilizing retrieved items is non-trivial. Applying Chain-of-Note \citep{yu2023chain} and structured data format \citep{yin-etal-2023-read} improves question answering accuracy by as much as 10 absolute points across three LLMs.
\end{itemize}
